\documentclass[12pt, a4paper]{article}

% --- Preambolo e Pacchetti ---
\usepackage[utf8]{inputenc}
\usepackage[italian]{babel}
\usepackage{amsmath}
\usepackage{amssymb}
\usepackage{amsfonts}
\usepackage{geometry}
\usepackage{tcolorbox}
\usepackage{enumitem}
\usepackage{hyperref}
\usepackage{bm}

% Configurazione margini
\geometry{margin=2.5cm}

% Configurazione tcolorbox per le formule
\tcbset{
    colback=blue!5!white,
    colframe=blue!75!black,
    fonttitle=\bfseries,
    center title,
    arc=3pt
}

\title{\textbf{Tecniche Quantistiche per la Sicurezza}\\ \large Lezione 1: Introduzione al Quantum Computing}
\author{Appunti del Corso}
\date{}

\begin{document}

\maketitle

\section{Introduzione al Corso}
Il corso non è focalizzato sulla fisica teorica pura, ma si propone come un percorso di \textbf{informatica applicata}. L'obiettivo è comprendere come utilizzare gli algoritmi quantistici, il \textit{Quantum Machine Learning} (QML) e la comunicazione quantistica per risolvere problemi di \textbf{cybersecurity} (es. anomaly detection, scambio di chiavi quantistiche).

La computazione quantistica viene oggi studiata poiché ci troviamo sulla frontiera della ricerca: i dispositivi reali (come quelli IBM) sono ancora limitati (pochi qubit e presenza di rumore), ma la tecnologia è in rapida espansione e promette di superare i limiti fisici dei computer classici.

\section{Rappresentazione Classica dei Dati}
Per capire il salto verso il quantistico, è necessario richiamare come i dati vengono gestiti nelle macchine classiche.

\subsection{Il Bit e il Transistor}
L'informatica classica si basa sul \textbf{bit} (0 o 1). Fisicamente, questo è realizzato tramite il \textbf{transistor}, che funge da interruttore elettronico.
\begin{itemize}
    \item \textbf{Struttura:} Un sandwich di materiali di tipo N e tipo P (transistor NPN).
    \item \textbf{Regioni di funzionamento:}
    \begin{enumerate}
        \item \textbf{Cut-off (Interdizione):} Il transistor è "OFF" (Logica 0). Non passa corrente.
        \item \textbf{Saturazione:} Il transistor è "ON" (Logica 1). La corrente passa liberamente.
        \item \textbf{Regione Attiva:} Usata per l'amplificazione (non usata nella logica digitale binaria perché instabile).
    \end{enumerate}
\end{itemize}

\paragraph{Perché la tensione e non la corrente?}
Si preferisce rappresentare i bit tramite \textbf{livelli di tensione} anziché di corrente perché la tensione è più stabile, meno influenzata dal carico e offre una migliore tolleranza al rumore grazie alla definizione di "zone sicure" (es. vicino a 0V è 0, vicino a 1V è 1).

\subsection{Limiti della Legge di Moore}
La \textit{Legge di Moore} prevede il raddoppio del numero di transistor ogni due anni. Tuttavia, stiamo raggiungendo la scala atomica (pochi nanometri), dove sorgono problemi insormontabili per la fisica classica:
\begin{itemize}
    \item \textbf{Quantum Tunneling (Tunneling Quantistico):} Quando le barriere isolanti diventano troppo sottili, gli elettroni possono "passare attraverso" a causa della loro natura ondulatoria, causando errori e rendendo il processore inaffidabile.
    \item \textbf{Dissipazione di calore:} L'impacchettamento di miliardi di transistor genera temperature elevatissime.
    \item \textbf{Complessità computazionale:} Molti problemi (es. simulazione molecolare o fattorizzazione RSA) crescono esponenzialmente con la dimensione dell'input, diventando trattabili in tempi biblici per i computer classici.
\end{itemize}

\section{Fondamenti di Meccanica Quantistica}
La meccanica quantistica descrive la materia e l'energia a scale subatomiche, dove l'intuizione classica fallisce.

\subsection{Dualismo Onda-Particella}
Le entità quantistiche (elettroni, fotoni) non si comportano come piccole palle da biliardo, ma hanno una natura duale. Sono descritte da una \textbf{funzione d'onda} $\psi$.

\begin{tcolorbox}[title=Lunghezza d'onda di De Broglie]
\[
\lambda = \frac{h}{p} = \frac{h}{mv}
\]
Dove $h$ è la costante di Planck e $p$ è il momento (massa $\times$ velocità).
\end{tcolorbox}

Il quadrato dell'ampiezza della funzione d'onda $|\psi(x)|^2$ fornisce la \textbf{densità di probabilità} di trovare la particella in una determinata posizione.

\subsection{Principio di Indeterminazione di Heisenberg}
Esiste un limite fondamentale alla precisione con cui possiamo conoscere coppie di proprietà fisiche coniugate (come posizione $x$ e momento $p$).

\begin{tcolorbox}[title=Principio di Indeterminazione]
\[
\Delta x \Delta p \ge \frac{\hbar}{2}
\]
Dove $\hbar = h / 2\pi$.
\end{tcolorbox}

\textbf{Conseguenza:} Se cerchiamo di misurare la posizione di un elettrone con estrema precisione, ne alteriamo il momento in modo imprevedibile. La misurazione stessa "disturba" il sistema.

\section{Il Qubit}
Il \textbf{qubit} è l'unità fondamentale dell'informazione quantistica. A differenza del bit classico, può esistere in uno stato di \textbf{superposizione}.

\subsection{Rappresentazione Matematica}
Un qubit è rappresentato come un vettore unitario in uno spazio complesso a due dimensioni, chiamato \textbf{Spazio di Hilbert} ($\mathcal{H} \cong \mathbb{C}^2$).

\begin{tcolorbox}[title=Stato del Qubit]
\[
|\psi\rangle = \alpha|0\rangle + \beta|1\rangle
\]
Con il vincolo di normalizzazione: $|\alpha|^2 + |\beta|^2 = 1$.
\end{tcolorbox}

\begin{itemize}
    \item \textbf{Ket ($|\cdot\rangle$):} Notazione di Dirac per un vettore colonna.
    \item \textbf{Ampiezze di probabilità:} $\alpha$ e $\beta$ sono numeri complessi.
    \item \textbf{Misurazione:} Quando misuriamo il qubit nella base computazionale $\{|0\rangle, |1\rangle\}$, otteniamo il risultato $0$ con probabilità $|\alpha|^2$ e il risultato $1$ con probabilità $|\beta|^2$. Dopo la misura, lo stato \textbf{collassa} nel valore ottenuto.
\end{itemize}

\subsection{La Sfera di Bloch}
La \textbf{Sfera di Bloch} è una rappresentazione geometrica intuitiva dello stato di un qubit:
\begin{itemize}
    \item Il polo nord rappresenta lo stato $|0\rangle$.
    \item Il polo sud rappresenta lo stato $|1\rangle$.
    \item I punti sulla superficie rappresentano \textbf{stati puri} (superposizioni coerenti).
    \item I punti all'interno della sfera rappresentano \textbf{stati misti} (sistemi con rumore o incertezza).
\end{itemize}

\section{Notazione Bra-Ket (Dirac)}
La notazione Bra-Ket è il linguaggio standard della meccanica quantistica per rappresentare stati e operatori.

\begin{itemize}
    \item \textbf{Ket $|\psi\rangle$:} Vettore colonna $\begin{pmatrix} \alpha \\ \beta \end{pmatrix}$.
    \item \textbf{Bra $\langle\psi|$:} Vettore riga, ottenuto come il complesso coniugato trasposto del Ket: $(\alpha^*, \beta^*)$.
\end{itemize}

\subsection{Prodotto Interno (Inner Product)}
Rappresenta lo scalare che deriva dal prodotto di un Bra e un Ket. Viene usato per calcolare ampiezze di probabilità.
\begin{tcolorbox}
\[
\langle\phi|\psi\rangle = (\gamma^*, \delta^*) \begin{pmatrix} \alpha \\ \beta \end{pmatrix} = \gamma^*\alpha + \delta^*\beta
\]
\end{tcolorbox}
\textbf{Proprietà:} Se $|\psi\rangle$ è normalizzato, $\langle\psi|\psi\rangle = 1$. Se due stati sono ortogonali, $\langle\phi|\psi\rangle = 0$.

\subsection{Prodotto Esterno (Outer Product)}
Il prodotto $|\phi\rangle\langle\psi|$ produce una \textbf{matrice} (operatore lineare).
\begin{itemize}
    \item \textbf{Proiettori:} Gli operatori $|0\rangle\langle 0|$ e $|1\rangle\langle 1|$ proiettano uno stato sulle basi corrispondenti.
    \item \textbf{Relazione di completezza:} $I = |0\rangle\langle 0| + |1\rangle\langle 1|$.
\end{itemize}

\section{Basi e Trasformazioni}
Oltre alla \textbf{base computazionale} (asse Z), esistono infinite basi. Le più importanti sono:
\begin{enumerate}
    \item \textbf{Base di Hadamard (Asse X):} Composta dagli stati $|+\rangle$ e $|-\rangle$.
    \begin{tcolorbox}
    \[
    |+\rangle = \frac{|0\rangle + |1\rangle}{\sqrt{2}}, \quad |-\rangle = \frac{|0\rangle - |1\rangle}{\sqrt{2}}
    \]
    \end{tcolorbox}
    \item \textbf{Base Circolare (Asse Y):} Composta dagli stati $|+i\rangle$ e $|-i\rangle$.
    \begin{tcolorbox}
    \[
    |+i\rangle = \frac{|0\rangle + i|1\rangle}{\sqrt{2}}, \quad |-i\rangle = \frac{|0\rangle - i|1\rangle}{\sqrt{2}}
    \]
    \end{tcolorbox}
\end{enumerate}

Cambiare base significa cambiare il "punto di vista" della misurazione. Misurare uno stato $|0\rangle$ nella base X darà risultati casuali (50\% $|+\rangle$ e 50\% $|-\rangle$).

\section{Teorema di Non-Clonazione}
\textbf{Definizione:} È fisicamente impossibile creare una copia identica e perfetta di uno stato quantistico arbitrario e sconosciuto.
Questa è una differenza fondamentale con l'informatica classica, dove il comando `copy` è banale. Ciò impedisce l'uso della ridondanza classica per la correzione degli errori e rende la comunicazione quantistica intrinsecamente sicura.

\section{Entanglement (Correlazione Quantistica)}
L'entanglement si verifica quando due o più qubit sono legati in modo tale che lo stato di uno non può essere descritto indipendentemente dall'altro, anche se separati da grandi distanze.

\subsection{Stati di Bell}
Sono i quattro stati di massimo entanglement per un sistema a due qubit.
\begin{tcolorbox}[title=Esempio: Stato di Bell $\Phi^+$]
\[
|\Phi^+\rangle = \frac{|00\rangle + |11\rangle}{\sqrt{2}}
\]
\end{tcolorbox}
In questo stato, se misuro il primo qubit e ottengo 0, so con \textbf{certezza istantanea} che anche il secondo qubit è 0, indipendentemente dalla distanza.

\paragraph{Paradosso EPR:} Einstein, Podolsky e Rosen misero in dubbio questa "azione a distanza", suggerendo l'esistenza di "variabili nascoste". Tuttavia, il \textbf{Teorema di Bell} e gli esperimenti successivi hanno dimostrato che la natura è intrinsecamente non-locale e quantistica.

\section{Rumore e Stati Misti}
Nel mondo reale, l'interazione con l'ambiente introduce rumore (\textit{Noisy World}). Lo stato non è più "puro" ma diventa un \textbf{Mixed State} (stato misto).

\subsection{Matrice di Densità}
Per descrivere uno stato misto (una distribuzione statistica di diversi stati puri $|\psi_i\rangle$ con probabilità $p_i$) si usa l'operatore di densità $\rho$.
\begin{tcolorbox}
\[
\rho = \sum_i p_i |\psi_i\rangle\langle\psi_i|
\]
\end{tcolorbox}
\textbf{Proprietà:} La traccia della matrice deve essere unitaria ($Tr(\rho) = 1$).

\subsection{Fidelity}
La \textbf{Fidelity} ($F$) misura quanto lo stato ottenuto $\rho$ è vicino allo stato ideale desiderato $|\psi\rangle$.
\begin{tcolorbox}
\[
F(\rho, |\psi\rangle) = \langle\psi|\rho|\psi\rangle
\]
\end{tcolorbox}
$F=1$ indica coincidenza perfetta; $F=0$ indica ortogonalità.

\section{Computazione Quantistica}
Una computazione quantistica consiste in tre fasi:
\begin{enumerate}
    \item \textbf{Preparazione (Encoding):} Si mappano i dati classici in stati quantistici (solitamente partendo da $|0\rangle$).
    \item \textbf{Trasformazione Unitaria:} Si applicano una serie di \textbf{gate quantistici} (operazioni reversibili).
    \item \textbf{Misurazione:} Si estrae l'informazione classica (operazione irreversibile e probabilistica).
\end{enumerate}

\subsection{Gate Quantistici Principali}
\begin{itemize}
    \item \textbf{Gate di Pauli:}
    \begin{itemize}
        \item \textbf{X (Bit-flip):} Inverte $|0\rangle \leftrightarrow |1\rangle$. Corrisponde al NOT classico.
        \item \textbf{Z (Phase-flip):} Cambia il segno della fase dello stato $|1\rangle$.
        \item \textbf{Y:} Combinazione di bit-flip e phase-flip.
    \end{itemize}
    \item \textbf{Gate di Hadamard (H):} Crea superposizione. Trasforma $|0\rangle \to |+\rangle$ e $|1\rangle \to |-\rangle$.
    \item \textbf{CNOT (Controlled-NOT):} Gate a due qubit. Inverte il secondo qubit (target) solo se il primo (control) è $|1\rangle$. Fondamentale per creare entanglement.
    \item \textbf{Toffoli (CCNOT):} Gate a tre qubit (due controlli, un target). È un gate universale per la logica classica.
\end{itemize}

\subsection{Insieme Universale (Universal Set)}
Un insieme di gate è universale se può approssimare con precisione arbitraria qualsiasi trasformazione unitaria. Un esempio comune è l'insieme $\{H, T, CNOT\}$. I gate \textbf{Clifford} (come $H, S, CNOT$) sono importanti per la correzione degli errori ma non bastano da soli per la computazione universale; serve almeno un gate non-Clifford come il gate $T$ ($\pi/8$).

\paragraph{Nota Operativa:} In laboratorio si utilizzeranno librerie come \textbf{Qiskit} (IBM) o \textbf{PennyLane}. Il processo di traduzione di un circuito logico in istruzioni per l'hardware fisico è chiamato \textbf{transpiling}.

\end{document}
